% STATUS: written 2026-11-05 (Bresar 1993 read at its first page, which
% carries the level definitions and Theorem A in full; Lee 1997 read in
% full; excerpts in excerpts/. Posner 1957 and Mayne are named through these
% later papers.)
% STATUS: extended 2026-11-06 (Franca and Louza 2024 read in full from the
% published PDF; the bars on P and E in Theorem 5.1 checked by word height).
% Additions of 2026-11-07 moved from chapter 20: Aczel 1981 (pp. 191--200) read in full.
% STATUS: extended 2026-11-06 (Tang 2016 (both papers), Fan, Su and Xu 2016, Pandey 2024 read in full)
\section{Commuting and centralizing maps}\label{sec:functional-identities:commuting-maps}

The question of this chapter is what a map must be when it behaves like a
derivation on products. A derivation obeys $d(xy)=d(x)y+xd(y)$, and the
commuting condition $[x,d(x)]=0$ is the special case that asks when a
derivation can agree with the element it acts on. This section treats maps
whose commutator with each element is central or zero, without assuming
additivity or a Leibniz rule.

A map $F$ of a ring $R$ into itself is \emph{commuting} on a subset $S$
\cite[p.~333]{LeeTK1997Centralizing} when $[s,F(s)]=0$ for all $s$ in $S$.
The map is \emph{centralizing} on $S$
\cite[p.~333]{LeeTK1997Centralizing} when $[s,F(s)]$ lies in the center
$Z$ of $R$.
The line began with a result of Posner from 1957: the existence of a
nonzero centralizing derivation on a prime ring forces the ring to be
commutative \cite[p.~385]{Bresar1993}. Mayne proved the
analogous statement for centralizing automorphisms
\cite[p.~385]{Bresar1993}. The
following years brought extensions of these two theorems in several
directions \cite[p.~385]{Bresar1993}.

Bre\v{s}ar asked for the structure of an arbitrary additive centralizing
map, and gave a complete answer in the prime case
\cite[p.~385]{Bresar1993}.

\begin{theorem}[{\cite[Theorem~A, p.~385]{Bresar1993}}]\label{thm:fi:bresarA}
Let $R$ be a prime ring. Suppose an additive mapping $F$ of $R$ into
itself is centralizing on $R$. If either $R$ has a characteristic
different from $2$ or $F$ is commuting on $R$, then $F$ is of the form
$F(x)=\lambda x+\zeta(x)$ for all $x\in R$, where $\lambda$ is an element
from the extended centroid $C$ of $R$ and $\zeta$ is an additive mapping
of $R$ into $C$.
\end{theorem}

Theorem~\ref{thm:fi:bresarA} says that a centralizing additive map on a
prime ring is, up to a central remainder, multiplication by a central
element \cite[p.~385]{Bresar1993}. Bre\v{s}ar later obtained the same
conclusion in the semiprime case, and based on these studies he initiated
the study of functional identities \cite[p.~333]{LeeTK1997Centralizing}.
For a right ideal the same form forces a restriction: every commuting
additive map $f\colon\rho\to R$ has the form $f(x)=\lambda x+\mu(x)$
unless $[\rho,\rho]\rho=0$ \cite[p.~333]{LeeTK1997Centralizing}.

The derivations between two centralizing maps are also controlled, at the
cost of restricting the ring \cite[p.~334]{LeeTK1997Centralizing}. A ring
is \emph{$S_4$-free} when it does not
satisfy the standard identity $S_4$ of degree $4$. Lee proved a result on
pairs of derivations \cite[Theorem~1, p.~334]{LeeTK1997Centralizing}.
Let $\delta$ and $d$ be nonzero derivations of a prime $S_4$-free ring with
$\delta\notin Cd$, and let $[\delta(x),d(x)]$ be central on a right ideal
$\rho$. Then $d$ and $\delta$ are inner on $\rho$ through two elements of
the Martindale quotient ring
\cite[Theorem~1, p.~334]{LeeTK1997Centralizing}.
The clean form of the centralizing theorem itself holds on one-sided
ideals \cite[Theorem~2, p.~335]{LeeTK1997Centralizing}. It moves
Theorem~\ref{thm:fi:bresarA} from the whole prime ring to a right ideal
\cite[p.~335]{LeeTK1997Centralizing}.

\begin{theorem}[{\cite[Theorem~2, p.~335]{LeeTK1997Centralizing}}]\label{thm:fi:leetheorem2}
Let $R$ be a prime ring and $\rho$ a nonzero right ideal of $R$. Then
every centralizing additive mapping $f\colon\rho\to U$ is of the form
$f(x)=\lambda x+\mu(x)$ for all $x\in\rho$, where $\lambda\in C$ and
$\mu\colon\rho\to C$, unless $[\rho,\rho]\rho=0$.
\end{theorem}

Here $U$ is the right Utumi quotient ring of $R$
\cite[p.~334]{LeeTK1997Centralizing}. Two earlier theorems sit behind
Theorem~\ref{thm:fi:leetheorem2} \cite[p.~334]{LeeTK1997Centralizing}.
An additive map $f$ of a prime ring with $[f(x),x]_n=0$ for a
fixed $n$ is already commuting, provided the characteristic is $0$ or
greater than $n$ \cite[Theorem~B, p.~334]{LeeTK1997Centralizing}. For a
derivation $d$ of a semiprime ring with $[d(x^n),x^n]_k=0$ on a right
ideal, the conclusion is $[\rho,R]d(R)=0$
\cite[Theorem~L, p.~334]{LeeTK1997Centralizing}. The general conjecture
that a positive-power condition $[f(x),x]_n=0$ already gives the form
$f(x)=\lambda x+\mu(x)$ remained open
\cite[p.~334]{LeeTK1997Centralizing}.

That conjecture is open, and this section states it as a problem. % claim-ok: own signpost introducing the moved problem, which carries its own citation

\begin{problem}[{\cite[p.~334]{LeeTK1997Centralizing}}]\label{prob:open:commuting}
Let $R$ be a semiprime ring and let $f\colon R\to R$ be an additive map with
$[f(x),x]_n=0$ for one fixed positive power $n$.
Must $f$ be of the form $f(x)=\lambda x+\mu(x)$, with $\lambda$ in the extended
centroid and $\mu$ a map into the extended centroid?
\end{problem}

Lee proves the statement for $n=1$ and for several restricted conditions on the
ring \cite[Theorem~L, p.~334]{LeeTK1997Centralizing}, and he records the
general power as open \cite[p.~334]{LeeTK1997Centralizing}.

So far the sets were a whole ring or a right ideal. Franca and Louza moved the
question to subsets of a matrix algebra studied in linear preserver problems, and
those subsets are not closed under addition \cite[p.~2540]{FrancaLouza2024}.

Throughout their paper $F$ is a field whose characteristic is $0$ or at least $5$, and
$M_n$ is the algebra of all $n\times n$ matrices over $F$ with $n\ge3$
\cite[p.~2541]{FrancaLouza2024}. Here $L$ is the prime field of $F$
\cite[p.~2541]{FrancaLouza2024}. The set $E$ consists of the idempotent matrices, and
$V$ consists of the scalar multiples $\alpha E_{ii}$ of the matrix units
\cite[p.~2541]{FrancaLouza2024}. Their sets are the union
$\overline{E}=E\cup V$, the set $S$ of matrices whose minimal polynomial has degree
$2$, and $P=\{X\in M_n\mid\operatorname{tr}(X)\in L\}$ with
$\overline{P}=P\cup V$ \cite[p.~2541]{FrancaLouza2024}.

An additive map $G$ on $M_n$ has the \emph{standard form} when
$G(X)=\lambda X+\mu(X)$ for all $X\in M_n$, where $\lambda$ is a scalar and $\mu$ is
an additive map into the scalar matrices \cite[p.~2540]{FrancaLouza2024}. In the
matrix algebra this is the conclusion of Theorem~\ref{thm:fi:bresarA}.

The first of the two answers goes to $S$ (Theorem~\ref{thm:fi:franca-louza-S}) \cite[p.~2550]{FrancaLouza2024}.

\begin{theorem}[{\cite[Theorem~6.6, p.~2550]{FrancaLouza2024}}]\label{thm:fi:franca-louza-S}
Let $F$ be a field whose characteristic is $0$ or at least $5$, let $n\ge3$, and let
$G\colon M_n\to M_n$ be an additive map which is commuting on $S$. Then $G$ has the
standard form.
\end{theorem}

So the condition on the matrices of minimal polynomial of degree $2$ is as strong as
the condition on the whole algebra. The idempotents behave in another way, and the
authors tie them to the trace condition \cite[p.~2547]{FrancaLouza2024}.

\begin{theorem}[{\cite[Theorem~5.1, p.~2547]{FrancaLouza2024}}]\label{thm:fi:franca-louza-E}
Let $F$ be a field whose characteristic is $0$ or at least $5$, let $n\ge3$, and let
$G\colon M_n\to M_n$ be an additive map. Then $G$ is commuting on $\overline{P}$ if
and only if $G$ is commuting on $\overline{E}$.
\end{theorem}

Theorem~\ref{thm:fi:franca-louza-E} says that the condition on the idempotents already
decides it on the larger set $\overline{P}$, which is still not the whole algebra
\cite[p.~2547]{FrancaLouza2024}. It does not force the standard form, and the authors
build a map that shows this \cite[Example~7.1, pp.~2550--2553]{FrancaLouza2024}.

\begin{example}[{\cite[Example~7.1, pp.~2550--2553]{FrancaLouza2024}}]\label{ex:fi:franca-louza}
Let $F$ be a field with $[F:L]\ge7$ and let $n=3$. Then there is an additive map
$G\colon M_3\to M_3$ which is commuting on $\overline{E}$ and does not have the
standard form. That map is not commuting on the set $R_1$ of rank-one matrices.
\end{example}

So $\overline{E}$ admits commuting maps outside the standard form, and $S$ does not
\cite[p.~2553]{FrancaLouza2024}. The authors compare the three classical sets in their
closing commentary: among $R_1$, $E$ and $S$ only $S$ has all its commuting maps of
the standard form \cite[p.~2553]{FrancaLouza2024}.

\section{Functional equations and information measures}\label{sec:functional-identities:information-functions}

A functional equation can carry a derivation inside its solutions
\cite[p.~195]{Aczel1981}. The presence of the derivation is what makes those
solutions surprising~\cite[p.~195]{Aczel1981}. Aczel studies one real function
of one variable, written $f$~\cite[p.~191]{Aczel1981}. The variable runs over
the open unit interval, and the function comes from a measure of
information~\cite[p.~191]{Aczel1981}. The measure of uncertainty behind it
satisfies a recursion~\cite[pp.~192--193]{Aczel1981}. Once the measure is
symmetric, the recursion becomes an equation on $f$ alone
\[
  f(x) + (1-x)\,f\!\left(\frac{y}{1-x}\right)
  = f(y) + (1-y)\,f\!\left(\frac{x}{1-y}\right),
\]
and $x$ and $y$ are positive with $x + y < 1$~\cite[pp.~192--193]{Aczel1981}.
Aczel calls the solutions the \emph{symmetric information
functions}~\cite[p.~192]{Aczel1981}.

The general solution of that equation involves a
derivation~\cite[p.~195]{Aczel1981}. Aczel writes the solution through two
auxiliary functions~\cite[p.~195]{Aczel1981}. The first one, written $e$, is
any solution of the product equation
\[
  e(xy) = x\,e(y) + y\,e(x), \qquad x, y \in {]0,1[} .
\]
The second one, written $d$, is any real derivation
\cite[pp.~191 and 194]{Aczel1981}:
\[
  d(x+y) = d(x) + d(y), \qquad d(xy) = x\,d(y) + y\,d(x).
\]
With those two in hand, the general solution reads
\[
  f(x) = e(x) + e(1-x) + ax
\]
for a constant $a$~\cite[p.~195]{Aczel1981}.

The point Aczel makes is about a regularity
assumption~\cite[p.~195]{Aczel1981}. The formula above is
correct~\cite[p.~195]{Aczel1981}. What fails is the claim that a regular $f$
forces $e$ to be regular in the same
class~\cite[p.~195]{Aczel1981}. The derivation is the
reason~\cite[p.~195]{Aczel1981}. If $e$ solves the two equations, then so does
$e + d$~\cite[p.~195]{Aczel1981}. A derivation satisfies both by definition,
and $d(1-x) = -d(x)$ follows from them~\cite[p.~195]{Aczel1981}. So a very
regular $f$ can hide a very irregular $e$~ (Theorem~\ref{thm:fi:aczel-general-solution})\cite[p.~195]{Aczel1981}.

\begin{theorem}[{\cite[p.~195]{Aczel1981}}]\label{thm:fi:aczel-general-solution}
Let $e_0$ be one particular solution of the two equations above for a given
$f$ and $a$. Then their general solution is $e = e_0 + d$, where $d$ is an
arbitrary real derivation, whether zero or not, restricted to ${]0,1[}$.
\end{theorem}

Aczel gives the converse argument too~\cite[p.~195]{Aczel1981}. Take two
solutions of the pair of equations~\cite[p.~195]{Aczel1981}. Their difference
satisfies both halves of the Leibniz
rule~\cite[p.~195]{Aczel1981}. So the difference is a
derivation~\cite[p.~195]{Aczel1981}.

The same effect appears again one degree
higher~\cite[p.~199]{Aczel1981}. Aczel also treats information functions of
degree $\alpha$~\cite[p.~199]{Aczel1981}. With no regularity assumption their
general solution is
\[
  f(x) = b\,[x^\alpha + (1-x)^\alpha - 1], \qquad \alpha \neq 1,
\]
so the case $\alpha = 2$ looked like the
rest~\cite[pp.~198--199]{Aczel1981}. The expectation was that the same
implication holds at $\alpha = 2$~\cite[p.~199]{Aczel1981}. Aczel reports a
counterexample instead~\cite[p.~199]{Aczel1981}.

\begin{example}[{\cite[p.~199]{Aczel1981}}]\label{ex:fi:aczel-alpha-two}
Let $d$ be a not identically zero real derivation and $b$ a constant. Then
\[
  f(x) = b\,[x^2 + (1-x)^2 - 1] + d(x), \qquad x \in {]0,1[},
\]
is a solution of the equation for degree $2$. It satisfies neither the
symmetry $f(x) = f(1-x)$ nor the equation for a degree other than $2$.
\end{example}

The derivation is what breaks the symmetry
here~\cite[p.~199]{Aczel1981}. A derivation changes sign under $x \mapsto
1-x$~\cite[p.~199]{Aczel1981}. The square bracket does not change sign under
that map~\cite[p.~199]{Aczel1981}. So the sum of the two cannot be symmetric.
Aczel adds a second observation~\cite[p.~199]{Aczel1981}. If an $f$ of this
shape is measurable, then the derivation in it is
zero~\cite[p.~199]{Aczel1981}. That is why the measurable solutions never
showed the effect~\cite[p.~199]{Aczel1981}.

So the derivations of this book are not always
constructive~\cite[p.~191]{Aczel1981}. Here they are the hidden term. They
separate the algebraic solutions of a functional equation from the regular
ones~\cite[p.~199]{Aczel1981}. That separation is what the two surprises in
the title of Aczel count~\cite[p.~191]{Aczel1981}.
The same question has an answer on infinite-dimensional Lie algebras
\cite[p.~1]{Tang2016W22}. On the Virasoro algebra every linear commuting map is $x\mapsto\lambda x+\sigma(x)c$,
with $\lambda$ a complex number and $\sigma$ a linear functional
\cite[Theorem~4.2, p.~9]{Tang2016W22}. On the $W$-algebra $W(2,2)$ a second term
appears, built from the map $\omega$ that replaces each generator $L_m$ by
$H_m$ \cite[Theorem~4.1, p.~9]{Tang2016W22}. The same two algebras carry no
non-trivial commutative post-Lie structure \cite[Theorem~4.6, p.~9]{Tang2016W22}.
On the deformed Schr\"odinger--Virasoro algebra the commuting maps have four
possible forms, according to the two parameters
\cite[Theorem~4.3, p.~11]{FanSuXu2016}. Each form is read from the
skew-symmetric biderivations of the algebra
\cite[Theorem~3.2, p.~4]{FanSuXu2016}.
For a Hom-Lie superalgebra the commuting condition carries graded signs, and a
linear super-commuting map on a centerless algebra equal to its derived algebra
lies in the centroid \cite[Theorem~4.2, p.~11]{Pandey2024Superbiderivations}.
The same paper describes the super-biderivations of that algebra
\cite[Theorem~3.8, p.~7]{Pandey2024Superbiderivations}. So on these algebras the
commuting maps, the biderivations and the centroid are three views of one
object \cite[p.~1]{Pandey2024Superbiderivations}.
